\begin{description}
	\item[1行目] 行数をオブジェクト\texttt{n}に入れる。行数を数える関数が\texttt{nrow}です。
	\item[2行目] データセットの列数と同じ長さの，1だけを繰り返し入れたベクトルを作ります。繰り返しは\texttt{rep}関数を使います。
	\item[3行目] 平均ベクトル$\bm{m}$を作る操作です。数式\ref{meanM}に対応する操作です。
	\item[4行目] 平均ベクトルを使って平均偏差行列$\bm{V}$を作る操作です。数式\ref{diffV}に対応しています。
	\item[5行目] 分散共分散行列$\bm{S}$を作る操作です。数式\ref{covS}に対応しています。
	\item[6行目] \texttt{diag}関数を使って，行列$\bm{S}$の対角要素を抜き出し，その平方根を取り出しています。\texttt{SD}は標準偏差が入ったベクトルになります。
	\item[7行目] ベクトルを\texttt{diga}関数にいれると，ベクトルの要素を対角に持つ正方行列ができます。それを$\bm{Q}$として表現しています。
	\item[8行目] 標準スコア行列$\bm{Z}$を作る操作です。平均偏差行列$\bm{V}$に$\bm{Q^{-1}}$をかけています。数式\ref{stdZ}に対応しています。
	\item[9行目] 相関行列$\bm{R}$を作る操作です。数式\ref{corR}に対応しています。
	\item[10行目] できあがった行列$\bm{R}$を検算するために，Rの持っている相関行列を作る関数\texttt{cor}を実行しています。作った$\bm{R}$と同じになっていることを確認してください。
\end{description}
